\begin{statement}{E.61}{Proved}{The exact mod-twelve genus-one sewing projection}
Use the positive generators of E.1 and E.58.  Let $(1,f)$ be any
normalized full gentle pair, with its actual extension
\[
 \Phi(s)=sf=\beta s\beta^{-1},\qquad
 \Phi(r)=f^{-1}r,\qquad \Phi(sr)=sr,
 \qquad \beta\in\widehat F.
 \tag{E.61.1}
\]
The pure frame $\beta$ is unique.  Put $h=\beta^{-1}$.  There is a
unique $\rho\in\widehat{\mathbb Z}$ with
\[
 h^{\rm ab}=(\rho,\rho),\qquad
 f(x,y)=h(y,x)^{-1}h(x,y).
 \tag{E.61.2}
\]
Thus $\rho$ is the Nakamura--Schneps frame coordinate in precisely
these generators.  Define the signed homomorphism
\[
 q_{12}:\widehat F\longrightarrow\mathrm{SL}_2(\mathbb Z/12),
 \qquad
 x\longmapsto\begin{pmatrix}1&2\\0&1\end{pmatrix},\quad
 y\longmapsto\begin{pmatrix}1&0\\-2&1\end{pmatrix}.
 \tag{E.61.3}
\]
Then every such pair satisfies the exact finite identity
\[
 q_{12}(f)=
 \begin{pmatrix}1&-8\rho\\-8\rho&1+64\rho^2\end{pmatrix}
                      \pmod {12}.
 \tag{E.61.4}
\]

Here is its direct genus-one comparison.  In $B_4=\langle a,b,c\rangle$,
put
\[
 E_1=(aba)^4,\qquad E_3=(cbc)^4,\qquad
 E_2=cba b^{-1}c^{-1}.
 \tag{E.61.5}
\]
The multiplier-one sewing word, with every frame term retained, is
\[
\begin{split}
 \mathcal R_\rho(f)={}&f(E_3,a)c^{-8\rho}f(b^2,c^2)
       f(E_2,E_3)f(E_1,E_2)\\
       &\hspace{13mm}\cdot f(a^2,b^2)a^{8\rho}f(c,E_1)
       \quad\hbox{in }\widehat B_4.
\end{split}
 \tag{E.61.6}
\]
For each integer $m\geq1$, the following is an actual continuous
homomorphism, where $(v,M)(w,N)=(v+Mw,MN)$:
\[
\begin{gathered}
 \eta_m:\widehat B_4\longrightarrow
       (\mathbb Z/m)^2\rtimes\mathrm{SL}_2(\mathbb Z/m),\\
 a\longmapsto(0,U),\qquad b\longmapsto(0,L),\qquad
 c\longmapsto(e_1,U),\qquad
 U=\begin{pmatrix}1&1\\0&1\end{pmatrix},\quad
 L=\begin{pmatrix}1&0\\-1&1\end{pmatrix}.
\end{gathered}
 \tag{E.61.7}
\]
Write $q_m(f)=f(U^2,L^2)=\left(\begin{smallmatrix}A&B\\C&D\end{smallmatrix}\right)$.
Then
\[
 \eta_m(\mathcal R_\rho(f))=
       \left(\binom{-B-8\rho A}{A-1},I\right).
 \tag{E.61.8}
\]
In particular $\eta_{12}(\mathcal R_\rho(f))=1$ for every normalized
full gentle pair of multiplier one.  The identity pair and the extra
CM pair of E.58 both have this value.  Thus this specified sewing
projection supplies no additional CM parity restriction.
\end{statement}
\begin{proof}
We first identify the frame and prove the finite matrix assertion;
we then compute the actual four-strand word.

Conjugation by $s$ interchanges $x$ and $y$.  The two expressions
for $\Phi(s)$ in (E.61.1) therefore give
$f=\beta(y,x)\beta(x,y)^{-1}$.
The centralizer calculation in E.58 gives
$C_{\widehat B}(s)\cap\widehat F=1$, so the pure frame is unique.
Write $\beta^{\rm ab}=(u,v)$.  Since $f\in\widehat F'$,
its abelianization $(v-u,u-v)$ is zero.  Hence $u=v$ and
$\rho=-u$ proves (E.61.2).  This is the word and abelianization
used in \cite[Definition 5.1]{BCNakamuraSchneps2000}; it is an
identity of the actual frames, with no comparison to a differently
normalized Kummer extension.

The image of $F$ modulo four under (E.61.3) is the abelian group
\[
 V_4=\left\{
 I,\begin{pmatrix}1&2\\0&1\end{pmatrix},
 \begin{pmatrix}1&0\\2&1\end{pmatrix},
 \begin{pmatrix}1&2\\2&1\end{pmatrix}\right\}.
 \tag{E.61.9}
\]
Thus $q_4(f)=I$.  To keep the sign at three, use the literal
matrices
\[
 S=\begin{pmatrix}0&1\\-1&0\end{pmatrix},\qquad
 R=\begin{pmatrix}0&-1\\1&-1\end{pmatrix}.
 \tag{E.61.10}
\]
They define $\widehat B\to\mathrm{SL}_2(\mathbb Z/12)/\{\pm I\}$,
and their restriction to $F$ is the projectivization of (E.61.3).
Indeed $SR=U^{-1}$, $x=(sr)^{-2}$, and $y=sxs$.
The two possibilities for the global sign are distinguished modulo
four, since $-I\notin V_4$.

Put $g=q_{12}(\beta)$.  The signed identity
\[
             q_{12}(f)=SgS^{-1}g^{-1}
 \tag{E.61.11}
\]
follows from the free-generator substitution, because $S$ conjugates
$U^2$ to $L^2$ and conversely.  Since $\Phi(r)^3=1$, the matrix
$(q_{12}(f)^{-1}R)^3$ is $I$ or $-I$.  Its reduction modulo four
is $R^3=I$, so it is $I$.  This proves the signed cubic condition
at three, without discarding the central sign.

For completeness, the required order-twenty-four calculation has
the following exact form.  Work in $\mathrm{SL}_2(\mathbb F_3)$,
and put $X=U^2$.  The character
\[
 \epsilon:\mathrm{SL}_2(\mathbb F_3)\longrightarrow\mathbb F_3,
 \qquad \epsilon(X)=\epsilon(L^2)=1
 \tag{E.61.12}
\]
has quaternion kernel.  Its existence and the complete enumeration
can also be checked by the decomposition
$\mathrm{SL}_2(\mathbb F_3)=\langle X\rangle Q_8$, where
$Q_8=\langle S,\left(\begin{smallmatrix}1&1\\1&2\end{smallmatrix}\right)\rangle$.
For $g$ in this group,
\[
 \bigl((SgS^{-1}g^{-1})^{-1}R\bigr)^3=I
 \quad\Longleftrightarrow\quad
 g=X^kS^j\quad(0\leq k<3,\ 0\leq j<4).
 \tag{E.61.13}
\]
The remaining twelve matrices give $-I$ for the cube.  Direct
multiplication on the twelve displayed matrices gives
\[
 SgS^{-1}g^{-1}
       =\begin{pmatrix}1&k\\ k&1+k^2\end{pmatrix},
       \qquad \epsilon(g)=k.
 \tag{E.61.14}
\]
These identities check every element, not merely the saved CM
representatives.  Now $\beta^{\rm ab}=(-\rho,-\rho)$, so
$\epsilon(g)=-2\rho=\rho$ modulo three.  Equations
(E.61.11)--(E.61.14) prove (E.61.4) modulo three.  Equation
(E.61.9) proves it modulo four, and the Chinese remainder
isomorphism proves the signed modulus-twelve assertion.

The three matrices in (E.61.7) satisfy the two adjacent braid
relations and the far commutation relation.  They therefore give
the indicated homomorphism.  Its linear part is the literal
Ferrari fold $B_4\to B_3$, $a,c\mapsto a$, $b\mapsto b$, followed
by the standard integral matrix representation.  The Ferrari
generators $ca^{-1}$ and
$aba c^{-1}b^{-1}a^{-1}$ map respectively to $e_1,e_2$.
Thus the translation module and its basis are fixed by actual
words.

Both $E_1$ and $E_3$ map to the identity under $\eta_m$.
For example the linear part of $cbc$ is $ULU$, whose square is
$-I$, and its four-term translation sum is zero.  Since
$f\in\widehat F'$, all four factors in (E.61.6) involving $E_1$
or $E_3$ therefore disappear.  We are left with
\[
 \eta_m(\mathcal R_\rho(f))
    =\eta_m(c)^{-8\rho}
       f(\eta_m(b)^2,\eta_m(c)^2)
       (0,q_m(f))\eta_m(a)^{8\rho}.
 \tag{E.61.15}
\]
Put $w=(0,-1)^t$ and $t_w=(w,I)$.  The two exact identities
\[
 t_w(0,L^2)t_w^{-1}=(0,L^2),\qquad
 t_w(0,U^2)t_w^{-1}=(2e_1,U^2)=\eta_m(c)^2
 \tag{E.61.16}
\]
prove, by the universal property of $\widehat F$, that
\[
 f(\eta_m(b)^2,\eta_m(c)^2)
       =\bigl((I-q_m(f)^{-1})w,q_m(f)^{-1}\bigr).
 \tag{E.61.17}
\]
Here $f(L^2,U^2)=q_m(f)^{-1}$ is precisely the gentle symmetry
identity.  The vector in (E.61.17) is $(-B,A-1)^t$.
Multiplying the four factors of (E.61.15) now gives
(E.61.8), including its signs and the factor eight.

We record the exact sewing-source comparison to prevent the
discarding of a central braid factor.  Write
\[
 X_{12}=a^2,\quad X_{23}=b^2,\quad X_{34}=c^2,\quad
 X_{45}=(aba)^2,\quad X_{51}=(cbc)^2,\quad Z=(abc)^4.
 \tag{E.61.18}
\]
Let $A_3=\langle u,v\mid uvuv=vuvu\rangle$ be the index-three
subgroup $\langle\tau_1,\tau_2^2\rangle$ of $B_3$.
The four actual homomorphisms to $B_4$ used by the sewing
calculation send $(u,v)$ to
\[
 (a,X_{51}),\quad(E_2,X_{51}),\quad
 (E_2,X_{45}),\quad(c,X_{45}).
 \tag{E.61.19}
\]
Their central squares $(uv)^2$ are, in order,
\[
          X_{34}Z,\quad X_{23}Z,\quad X_{23}Z,\quad X_{12}Z.
 \tag{E.61.20}
\]
These are equalities in the discrete braid group.  The certificate
checks the defining relation and each square through the faithful
Artin representation, whose generators are written explicitly in
the output; faithfulness in this convention is
\cite[Theorem A, representation (1) with $k=1$]{BCShpilrain2002}.
Thus the maps extend to profinite completions.

For a pair in the Nakamura--Schneps subgroup, relation (IV) of
\cite[Section 1]{BCNakamuraSchneps2000} reads
\[
 f(u,v^2)=v^{4\rho}f(u^2,v)u^{4\rho}(uv)^{-4\rho}.
 \tag{E.61.21}
\]
Applying the four maps, with inversion on the first and third
expressions, retains the following full braid formulas:
\[
\begin{aligned}
 f(E_3,a)&=Z^{2\rho}X_{34}^{2\rho}X_{12}^{-2\rho}
                    f(X_{51},X_{12})X_{51}^{-4\rho},\\
 f(E_2,E_3)&=Z^{-2\rho}X_{51}^{4\rho}f(X_{14},X_{51})
                    X_{14}^{2\rho}X_{23}^{-2\rho},\\
 f(E_1,E_2)&=Z^{2\rho}X_{23}^{2\rho}X_{14}^{-2\rho}
                    f(X_{45},X_{14})X_{45}^{-4\rho},\\
 f(c,E_1)&=Z^{-2\rho}X_{45}^{4\rho}f(X_{34},X_{45})
                    X_{34}^{2\rho}X_{12}^{-2\rho}.
\end{aligned}
 \tag{E.61.22}
\]
Here $X_{14}=E_2^2$.  The central powers in their product cancel
with coefficients $2,-2,2,-2$.  This is the literal multiplier-one
word of \cite[proof of Claim 8.3, relation $(R')$]{BCNakamuraSchneps2000},
which is also the genus-one sewing calculation invoked in
\cite[Theorems 3.18 and 4.9]{BCBasualdoRobertson2026}.
Equation (E.61.22) is the source's calculation on its
Nakamura--Schneps subgroup; the proof of (E.61.4) and
(E.61.8) did not impose (IV) on a full-GT element.

The program \texttt{bc\_genus\_one\_sewing.py} verifies all sixteen
displayed braid comparisons, the complete twenty-four-element
matrix classification, the translation conjugator, and the actual
framed words of E.1.  Its JSON retains the four-generator Artin
images, the twelve passing and twelve failing cubic matrices,
the exact recovered $\beta$ words and $\rho$ values, and every
mod-three, mod-four and mod-twelve sewing value.  All twenty-four
saved CM words give the identity in these projections.
The initial search for a change under the $4,656$ literal
$N_B'$ commutator corrections to $\beta$ found none; it is
preserved under \texttt{research/E61\_initial\_affine\_probe}.
The universal calculation above explains that outcome.
\end{proof}
